% Created 2024-01-08 lun 09:26
% Intended LaTeX compiler: pdflatex
\documentclass[11pt]{article}
\usepackage[utf8]{inputenc}
\usepackage[T1]{fontenc}
\usepackage{graphicx}
\usepackage{longtable}
\usepackage{wrapfig}
\usepackage{rotating}
\usepackage[normalem]{ulem}
\usepackage{amsmath}
\usepackage{amssymb}
\usepackage{capt-of}
\usepackage{hyperref}
\usepackage{minted}

\usepackage{parskip}
\usepackage[utf8]{inputenc} %% For unicode chars
\usepackage{placeins}
\usepackage[margin=2.50cm]{geometry}
\usepackage[T1]{fontenc}
\usepackage{mathpazo}
\linespread{1.05}
\usepackage[scaled]{helvet}
\usepackage{courier}
\hypersetup{colorlinks=true,linkcolor=blue}
\RequirePackage{fancyvrb}
\DefineVerbatimEnvironment{verbatim}{Verbatim}{fontsize=\small,formatcom = {\color[rgb]{0.5,0,0}}}
\usepackage{mdframed}
\BeforeBeginEnvironment{minted}{\begin{mdframed}}
\AfterEndEnvironment{minted}{\end{mdframed}}
\setcounter{secnumdepth}{3}
\date{}
\title{PHP Object-Oriented Programming (OOP)}
\hypersetup{
 pdfauthor={},
 pdftitle={PHP Object-Oriented Programming (OOP)},
 pdfkeywords={},
 pdfsubject={},
 pdfcreator={Emacs 28.1 (Org mode 9.5.5)}, 
 pdflang={Spanish}}
\begin{document}

\maketitle
\setcounter{tocdepth}{3}
\tableofcontents

\setlength\parindent{10pt}


\section{PHP Object-Oriented Programming (OOP)}
\label{sec:org6ed5851}
Object-oriented programming (OOP) is a programming paradigm that
organizes software using objects. Objects are self-contained units of
data that encapsulate both data and the operations that can be performed
on it. OOP provides several advantages over procedural programming,
including:

\begin{itemize}
\item \textbf{Modular design:} OOP promotes modular design by dividing programs
into independent objects that can be easily reused and composed.

\item \textbf{Encapsulation:} OOP facilitates data encapsulation, which means that
data is protected from direct access and can only be modified through
the object's methods. This promotes data integrity and prevents
accidental data corruption.

\item \textbf{Inheritance:} OOP enables code reuse by allowing new classes to
inherit properties and methods from existing classes. This reduces
code duplication and promotes code maintainability.

\item \textbf{Polymorphism:} OOP supports polymorphism, which allows objects of
different classes to respond to the same method call in different
ways. This promotes flexibility and adaptability in code.
\end{itemize}

PHP is a scripting language that supports OOP concepts. It provides a
rich set of features for defining classes, creating objects, and
managing object interactions. This guide provides a comprehensive
overview of OOP in PHP, covering essential concepts, advanced
techniques, and best practices.

\subsection{\textbf{Class}}
\label{sec:org8d7080a}
A class is the blueprint for creating objects. It defines the structure
of an object, including its properties (data) and methods (behaviors).

\textbf{Object Instantiation}

An object is an instance of a class. It represents a specific instance
of the data and behavior defined in the class. To create an object, use
the \texttt{new} keyword followed by the class name.

\begin{minted}[]{php}
<?=
  class Student {
      public $name;
      public $age;

      public function __construct($name, $age) {
          $this->name = $name;
          $this->age = $age;

          echo "Creating a new student: {$this->name} ({$this->age})";
      }
  }

  $student1 = new Student('John Doe', 25);
  $student2 = new Student('Jane Doe', 22);
\end{minted}

\textbf{Accessing Object Members}

Object members (properties and methods) can be accessed using the dot
notation. For properties, use the accessor \texttt{\$this->} to access them
within the object's methods. For methods, use the object name followed
by the method name and parentheses to call the method.

\begin{minted}[]{php}
<?=
$student1->name = 'John Smith'; // Set the name property
echo $student1->name; // Get the name property

$student2->sayHello(); // Call the sayHello() method
\end{minted}

\textbf{Initial Properties Values}

Properties can be assigned initial values when declaring the class.
These values are assigned to the object's properties when the object is
instantiated.

\begin{minted}[]{php}
<?=
class Student {
    public $name = 'John Doe';
    public $age = 25;
}

$student1 = new Student();
echo $student1->name; // Output: John Doe
echo $student1->age; // Output: 25
\end{minted}

\textbf{Constructor}

The constructor is a special method that is called when an object is
created. It is typically used to initialize the object's properties. The
constructor is defined using the \texttt{\_\_construct()} method.

\begin{minted}[]{php}
<?=
class Student {
    public function __construct($name, $age) {
        $this->name = $name;
        $this->age = $age;
    }
}

$student1 = new Student('John Doe', 25); // Initialize the object's properties
\end{minted}

\textbf{Destructor}

The destructor is a special method that is called when an object is
destroyed. It is typically used to perform cleanup tasks, such as
closing database connections or releasing resources.

\begin{minted}[]{php}
<?=
class Student {
    public function __destruct() {
        echo "Destroying student object";
    }
}

$student1 = new Student('John Doe', 25);
$student2 = new Student('Jane Doe', 22);

unset($student1); // Destroy the student1 object
 unset($student2); // Destroy the student2 object
\end{minted}

\textbf{Case Sensitivity}

PHP class names and object names are case-sensitive. This means that
\texttt{Student} and \texttt{student} represent different classes or objects.

\textbf{Object Comparison}

Objects can be compared using the \texttt{==} and \texttt{!=} operators. These
operators compare the references of objects, not their values. To
compare the values of objects, use the appropriate accessor methods or
properties. For example, to compare the name property of two Student
objects, you would use the following code:
\begin{minted}[]{php}
<?=
$student1 = new Student('John Doe', 25);
$student2 = new Student('Jane Doe', 22);

if ($student1->getName() == $student2->getName()) {
    echo "The students have the same name.";
} else {
    echo "The students do not have the same name.";
}
\end{minted}

If you want to compare objects based on their type, you can use the
\texttt{instanceof} operator. For example, to check if an object is an instance
of the Student class, you would use the following code:
\begin{minted}[]{php}
<?=
  $person = new Student('John Doe', 25);

  if ($person instanceof Student) {
      echo "The person is a student.";
  } else {
      echo "The person is not a student.";
  }
\end{minted}

\subsection{\textbf{Anonymous Class}}
\label{sec:org059b149}
An anonymous class, also known as a lambda or closure, is a class
without a name. It is defined using the \texttt{new} keyword followed by a
block of code. Anonymous classes are often used for creating temporary
objects or for passing objects as function arguments.

\begin{minted}[]{php}
$anonymousStudent = new class ('John Doe', 25) {
    public function sayHello() {
        echo "Hello, my name is {$this->name}.";
    }
};

$anonymousStudent->sayHello(); // Output: Hello, my name is John Doe.
\end{minted}

\subsection{\textbf{Inheritance}}
\label{sec:org5820731}
Inheritance is a key concept in OOP that allows new classes to inherit
properties and methods from existing classes. This promotes code reuse
and reduces code duplication.

\textbf{Overriding Members}

Classes that inherit from other classes can override existing properties
and methods. This allows for the customization of behavior without
modifying the original class.

\begin{minted}[]{php}
class Animal {
    public function makeSound() {
        echo "The animal makes a sound.";
    }
}

class Dog extends Animal {
    public function makeSound() {
        echo "Woof!";
    }
}

$dog = new Dog();
$dog->makeSound(); // Output: Woof!
\end{minted}

\textbf{Final Keyword}

The \texttt{final} keyword can be used to prevent a class or method from being
overridden. This ensures that the original implementation is not changed
by subclasses.

\begin{minted}[]{php}
class Animal {
    final public function makeSound() {
        echo "The animal makes a sound.";
    }
}

class Dog extends Animal {
    // Cannot override the final makeSound() method
    public function makeSound() {
        echo "Woof!";
    }
}
\end{minted}

\textbf{instanceof Operator}

The \texttt{instanceof} operator is used to check whether an object is an
instance of a particular class.

\begin{minted}[]{php}
$dog = new Dog();
if ($dog instanceof Animal) {
    echo "The dog is an instance of the Animal class.";
} else {
    echo "The dog is not an instance of the Animal class.";
}
\end{minted}

\subsection{\textbf{Access Levels}}
\label{sec:orgd996984}
Access levels determine the accessibility of properties and methods
within a class. There are three access levels in PHP:

\begin{itemize}
\item \textbf{Private:} Properties and methods declared as \texttt{private} can only be
accessed within the class itself.

\item \textbf{Protected:} Properties and methods declared as \texttt{protected} can be
accessed within the class and within subclasses of the class.

\item \textbf{Public:} Properties and methods declared as \texttt{public} can be accessed
anywhere in the program.
\end{itemize}

\textbf{var Keyword}

The \texttt{var} keyword is deprecated in PHP and is not recommended for use.
It is better to use explicit access modifiers for clear and concise
code.

\textbf{Object Access}

Access modifiers determine the access level of class members when
accessing them from outside the class. For example, a \texttt{public} property
can be accessed directly using the object name, while a \texttt{protected}
property needs to be accessed using the \texttt{->} operator and a \texttt{public}
method.

\begin{minted}[]{php}
class Animal {
    public $name;

    protected $age;

    public function getAge() {
        return $this->age;
    }
}

$dog = new Animal();
$dog->name = 'Max'; // Access public property directly
$dog->getAge(); // Access protected property through method
\end{minted}

\textbf{Access Level Guideline}

In general, it is recommended to use \texttt{private} for properties that
should not be accessed outside the class, \texttt{protected} for properties
that should only be accessible within the class and subclasses, and
\texttt{public} for properties and methods that need to be accessed from
anywhere in the program.

\textbf{Accessors}

Accessors are methods that provide controlled access to properties. They
are used to set and get property values while also performing validation
or other operations.

Sure, here is the completed code for the \texttt{Animal} class:

\begin{minted}[]{php}
class Animal {
    private $name;
    private $age;

    public function setName($name) {
        if (is_string($name)) {
            $this->name = $name;
        } else {
            throw new Exception('Invalid name');
        }
    }

    public function getName() {
        return $this->name;
    }

    public function setAge($age) {
        if (is_int($age) && $age >= 0) {
            $this->age = $age;
        } else {
            throw new Exception('Invalid age');
        }
    }

    public function getAge() {
        return $this->age;
    }
}
\end{minted}

The code includes additional validation for the \texttt{setAge()} method to
ensure that the age is an integer and a non-negative value. This helps
to maintain data integrity and prevent errors caused by invalid age
values.

\subsection{\textbf{Static}}
\label{sec:org3168627}
Static members are shared amongst all instances of a class and are not
bound to specific objects. They are declared using the \texttt{static} keyword.

\textbf{Referencings Static Members}

Static members can be accessed using the class name instead of an
object.

\begin{minted}[]{php}
class Animal {
    public static $species = 'Mammal';

    public function __construct() {
        static::$species = 'Dog'; // Modify static property
    }
}

$dog = new Animal();
echo Animal::$species; // Output: Dog
\end{minted}

\textbf{Static Variables}

Static variables can be used to store shared data that is independent of
object instances.

\begin{minted}[]{php}
class Animal {
    private static $count = 0;

    public function __construct() {
        static::$count++;
    }
}

$dog1 = new Animal();
$dog2 = new Animal();

echo Animal::$count; // Output: 2
\end{minted}

\textbf{Late Static Binding}

Late static binding allows method calls to be resolved at runtime, even
when the method is not yet defined. This is useful for dynamic class
loading and inheritance hierarchies.

\begin{minted}[]{php}
class Animal {
    public static function sayHello() {
        return 'The animal says hello.';
    }

    public static function getGreetingMethod() {
        return self::sayHello; // Late static binding
    }
}

class Dog extends Animal {
    public static function sayHello() {
        return 'Woof!';
    }
}

$dog = new Dog();
echo Animal::getGreetingMethod(); // Outputs: Woof!
\end{minted}

\subsection{\textbf{Interface}}
\label{sec:org7144e98}
An interface defines a set of abstract methods that must be implemented
by classes that implement the interface. Interfaces are used to enforce
contracts and promote code reusability.

\textbf{Interface Signatures}

Interfaces specify the method names, argument lists, and return types.
They do not contain implementation details.

\begin{minted}[]{php}
interface AnimalInterface {
    public function makeSound();
}

class Dog implements AnimalInterface {
    public function makeSound() {
        echo "Woof!";
    }
}

class Cat implements AnimalInterface {
    public function makeSound() {
        echo "Meow!";
    }
}
\end{minted}

\textbf{Interface Example}

Interfaces can be used to define common behavior for related classes.

\begin{minted}[]{php}
interface ShapeInterface {
    public function calculateArea();
}

class Square implements ShapeInterface {
    private $sideLength;

    public function __construct($sideLength) {
        $this->sideLength = $sideLength;
    }

    public function calculateArea() {
        return $this->sideLength * $this->sideLength;
    }
}

class Circle implements ShapeInterface {
    private $radius;

    public function __construct($radius) {
        $this->radius = $radius;
    }

    public function calculateArea() {
        return pi() * pow($this->radius, 2);
    }
}
\end{minted}

\textbf{Interface Usage}

Interfaces can be used to create objects that can interact with each
other based on their common behavior.

\begin{minted}[]{php}
$square = new Square(5);
echo $square->calculateArea(); // Output: 25

$circle = new Circle(3);
echo $circle->calculateArea(); // Output: 28.26
\end{minted}

\textbf{Interface Guideline}

Use interfaces to define common behavior for related classes and promote
code reusability.

\subsection{\textbf{Abstract}}
\label{sec:org0bebb5d}
An abstract class is a class that contains at least one abstract method.
Abstract methods do not have implementations and must be implemented by
subclasses. Abstract classes are used to define partial implementations
and enforce contracts.

\textbf{Abstract Methods}

Abstract methods are declared using the \texttt{abstract} keyword.

\begin{minted}[]{php}
abstract class AnimalAbstract {
    abstract public function makeSound();
}

class Dog extends AnimalAbstract {
    public function makeSound() {
        echo "Woof!";
    }
}

class Cat extends AnimalAbstract {
    public function makeSound() {
        echo "Meow!";
    }
}
\end{minted}

Abstract classes and interfaces are both powerful tools in
object-oriented programming (OOP) that can be used to define and
structure classes. While they share some similarities, they also have
distinct differences.

\textbf{Similarities:}

\begin{itemize}
\item \textbf{Encapsulation:} Both abstract classes and interfaces promote
encapsulation by encouraging the use of protected and private access
modifiers for properties and methods. This helps to protect data from
unauthorized access and maintain data integrity.

\item \textbf{Code Reusability:} Both abstract classes and interfaces can be used
to promote code reusability by defining common behavior or structure
that can be inherited by subclasses. This reduces redundancy and
improves maintainability of code.

\item \textbf{Type Safety:} Both abstract classes and interfaces can help to
improve type safety by specifying the types of data expected and
returned by methods. This helps to catch potential errors and maintain
code consistency.
\end{itemize}

\textbf{Differences:}

\begin{itemize}
\item \textbf{Purpose:} Abstract classes are more focused on defining the overall
structure or blueprint of a class, including its attributes and
behavior. They can contain both abstract and concrete methods.
Interfaces, on the other hand, are specifically designed to define a
set of behaviors that must be implemented by classes that implement
the interface. They only contain abstract methods.

\item \textbf{Implementation:} Abstract classes can contain concrete methods, which
provide default implementations for methods that are likely to be
similar across subclasses. Interfaces, however, do not contain any
implementation details, leaving that responsibility entirely to the
implementing classes.

\item \textbf{Relationship:} Abstract classes can be related to each other through
inheritance, forming a hierarchy of classes. Interfaces, on the other
hand, are more independent of each other, and a class can implement
multiple interfaces without being related to them through inheritance.
\end{itemize}

\textbf{When to Use:}

\begin{itemize}
\item \textbf{Abstract Classes:} Use abstract classes when you want to define a
common structure or blueprint for related classes, including both
common properties and behavior. You can also use abstract classes to
enforce encapsulation and promote type safety.

\item \textbf{Interfaces:} Use interfaces when you want to define a set of
behaviors that must be implemented by classes. Interfaces are
particularly useful for decoupling classes and promoting polymorphism.
\end{itemize}

In summary, abstract classes and interfaces are both valuable tools in
OOP, each with its unique strengths and applications. Abstract classes
are more focused on defining class structure, while interfaces are
specifically designed to define behaviors. The choice between the two
depends on the specific requirements and design of your application.

\subsection{\textbf{Traits}}
\label{sec:orgf949904}
Traits are a feature introduced in PHP 5.4 that allows the sharing of
functionality between unrelated classes. Traits can contain methods,
properties, and constants.

\textbf{Inheritance \& Traits}

Traits can be inherited from multiple classes, and their methods can be
overridden or extended.

\begin{minted}[]{php}
trait Logger {
    public function log($message) {
        echo "Logging: $message";
    }
}

class Animal {
    use Logger;

    public function makeSound() {
        echo "The animal makes a sound.";
    }
}

class Dog extends Animal {
    public function makeSound() {
        echo "Woof!";
        $this->log('Dog made a sound.');
    }
}

$dog = new Dog();
$dog->makeSound(); // Logs "Logging: Dog made a sound."
\end{minted}

\textbf{Trait Guidelines}

Use traits to share common functionality between unrelated classes
without the need for inheritance.

\subsection{\textbf{Overloading}}
\label{sec:orgc82881b}
Overloading is a technique in OOP that allows multiple methods with the
same name to exist in a class. Overloading is based on the method's
parameter types and order.

\textbf{Property Overloading}

Property overloading is not supported in PHP. However, it is possible to
achieve similar functionality using setters and getters.

\textbf{Method Overloading}

Method overloading is supported in PHP using the same method name and
different parameter lists.

\begin{minted}[]{php}
class Number {
    public function add($x, $y) {
        return $x + $y;
    }

    public function add($x, $y, $z) {
        return $x + $y + $z;
    }
}

$number = new Number();
echo $number->add(10, 20); // Output: 30
echo $number->add(10, 20, 30); // Output: 60
\end{minted}

\textbf{isset \& unset Overloading}

\texttt{isset()} and \texttt{unset()} can be overloaded to check and unset properties
based on their types.

\begin{minted}[]{php}
class Number {
    public function __isset($name) {
        if (is_int($name)) {
            return $this->$name !== null;
        } else {
            return false;
        }
    }

    public function __unset($name) {
        if (is_int($name)) {
            $this->$name = null;
        }
    }
}

$number = new Number();
$number->x = 10;

if (isset($number->x)) {
    echo "x is set";
} else {
    echo "x is not set";
} // Output: x is set

unset($number->x);

if (isset($number->x)) {
    echo "x is set";
} else {
    echo "x is not set";
} // Output: x is not set
\end{minted}

Overall, OOP is a powerful and versatile programming paradigm that can
significantly improve the maintainability, reusability, and flexibility
of PHP applications.
\end{document}
